\chapter{Component and market records}
\label{app:component-records}

Chapter~\ref{ch:software-dossiers} follows one research note through the
component decision. This appendix preserves the fuller technical and market
record behind that argument. A prose linter may catch a repeated opener and
miss a changed equation. A LaTeX parser may preserve structure and still say
nothing about whether a reader understands the request. Each candidate is
therefore recorded by the job it can perform, not the size of its feature list.

The inventory is a snapshot, not a promise about future releases. Vendored
repositories have pinned commits in the audit record; remote-only candidates
are identified as such. A documentation page establishes availability and
intended use. It does not establish effectiveness on Humanvoice documents.
The disposition of every candidate therefore remains conditional on a fixture
replay and a human-burden observation.

The adoption path in Figure~\ref{fig:package-funnel} prevents four different
kinds of evidence from collapsing into ``supported.'' A listed project has not
yet earned a runtime role, and a successful installation has not yet earned an
editor's attention.

\section{The capability map}
\label{sec:software-capabilities}

The first build has to do six jobs that vendors often sell as one ``AI
writing'' feature. An author needs to be able to:

\begin{enumerate}
\item parse a source document and retain locations for equations, tables,
      citations, labels, and prose;
\item build the document and expose cross-reference or compilation failures;
\item locate lexical, structural, or citation findings without changing text;
\item compare protected objects across two revisions;
\item run optional local inference while keeping source material inside the
      declared trust boundary; and
\item evaluate a rendered document with a named reader and a reproducible
      record.
\end{enumerate}

No candidate in the current inventory performs all six. That is not a defect in
the individual projects. It is a reason to compose small components behind a
narrow Humanvoice record instead of asking one package to become an editor,
parser, judge, and compliance system at once.

\begin{table}[H]
\centering
\small
\begin{tabularx}{\textwidth}{@{}p{0.19\textwidth}p{0.25\textwidth}L p{0.21\textwidth}@{}}
\toprule
Job & Candidate families & Evidence a candidate can provide & Evidence it
cannot provide alone \\
\midrule
Source representation & pylatexenc, plasTeX, LaTeXML, tree-sitter-LT, Pandoc &
Locations, syntax trees, conversion behavior on fixtures & Semantic correctness
of an equation or reader comprehension \\
Build and cross-reference & TeX, latexdiff, LaTeX tooling & Reproducible PDF,
undefined-reference and visual-diff observations & Preservation of an author's
intended claim \\
Prose findings & Vale, textlint, LanguageTool, proselint, Textidote & Rule hits,
locations, configurable style checks & Priority, genre fit, or a useful repair \\
AI-tell diagnostics & sloptrim, slop-forensics, fast-detect-gpt & Aggregate
distributional signals under a declared corpus & Individual authorship or
quality \\
Private inference & llama.cpp and compatible local models & A local execution
path and resource profile & Truth, citation support, or reader acceptance \\
Evaluation & Inspect AI plus a human protocol & Replayable tasks and judge
adapters & The named expert's final decision \\
\bottomrule
\end{tabularx}
\caption{The package field supplies components, not a finished product.}
\label{tab:software-capability-map}
\end{table}

\section{Representing and building technical documents}
\label{sec:software-parsers}

\subsection{Lightweight and full-model LaTeX parsers}

pylatexenc offers a relatively small Python-level route for tokenizing and
interpreting LaTeX fragments \citep{pylatexenc2026}. Its attraction is
proportionate scope: a first adapter can recover command spans, groups, and
source offsets without committing the product to a full document model. The
failure case is equally clear. A token stream does not know whether a macro
represents a theorem, a number in a table, or an author's private shorthand.
The adapter must retain the raw span and abstain when a macro cannot be
interpreted.

plasTeX builds a richer object model and can expose sectioning, counters, and
cross-references \citep{plastex2026}. It is a stronger candidate when the
product needs to reason about a document's visible hierarchy. It also brings a
larger compatibility surface: custom packages, macro definitions, and
environment side effects can make the object tree differ from the PDF a local
TeX installation produces. The benchmark must compare both the parsed model
and the compiled output.

LaTeXML translates TeX into XML and HTML-oriented structures and is useful when
the same source must support multiple renderings \citep{latexml2026}. Its
strength is explicit semantic markup for many mathematical constructs. Its
boundary is practical: a translation that looks correct in HTML can still
alter spacing, numbering, or a custom macro in the PDF used by the reader.
LaTeXML therefore belongs in an adapter experiment, not as an unexamined
replacement for the project's build.

unified-latex adds a serious JavaScript/TypeScript candidate. Its parser and
unified-style transformation utilities make it attractive for an editor or web
interface that needs source positions and typed nodes
\citep{unifiedlatex2026}. TexSoup sits at the opposite end of the design space:
its tolerant Python interface makes navigation and small transformations easy,
including on imperfect input \citep{texsoup2026}. Tolerance is useful for an
editor and dangerous for a safety claim. Both candidates must expose unknown
constructs and preserve raw spans; neither may convert a best-effort parse into
a false ``protected'' result.

\subsection{tree-sitter, Pandoc, and latexdiff}

tree-sitter's incremental parsing model is attractive for an editor that must
show a finding while a source file is being changed
\citep{treesitterlatex2026}. The current grammar is a structural aid, not a
semantic LaTeX proof. A malformed environment or an unsupported macro should
produce an explicit ``unparsed'' record, not a best-effort protected-object
comparison.

Pandoc supplies mature conversion and filter hooks across document formats
\citep{pandoc2026filters}. It can help create a detexed prose view for
aggregate diagnostics and a plain-text reader packet. It cannot be the source
of truth for equation preservation when the canonical source is LaTeX. The
source-to-PDF build remains authoritative; Pandoc output is a secondary view
with its own hash and conversion report.

latexdiff is valuable precisely because it is narrow. It highlights visible
source changes and helps a reader inspect a revision, but its markup is not a
semantic correspondence algorithm \citep{latexdiff2026}. Humanvoice should
use it as a view over the protected comparison, never as the comparison
itself. A number can move from a macro definition to a table, and an equation
can remain algebraically equivalent while its text changes; character markup
cannot resolve either case by itself.

The canonical build should also use the standard LaTeX toolchain rather than
hide it behind a bespoke command. latexmk can orchestrate repeat compilation
and bibliography dependencies \citep{latexmk2026}; ChkTeX can add source-level
warnings \citep{chktex2026}. These tools establish build and syntax facts. They
do not inspect rendered pages, compare protected claims, or decide whether a
reader understands the result, so they remain components behind the Humanvoice
build record rather than substitutes for it.

\begin{cardexample}{a parser failure that must remain visible}
The fixture defines a custom command \texttt{\\riskterm} that expands to a
number in one table and to a prose label in another. A parser that records only
the expanded text can report two unrelated values; a parser that preserves the
macro call and its expansion can ask the author to confirm the correspondence.
The second behavior is slower and safer. The first is an attractive false pass.
\end{cardexample}

\section{Prose and citation diagnostics}
\label{sec:software-prose}

Vale is a rule-based prose linter with configurable styles and targeted text
segments \citep{vale2026docs}. It is a good host for explicit project rules:
private path leakage, a banned internal phase label, or a required definition
of an acronym. The rule should report the passage and a reason. A rule that
rewrites the sentence or treats a hit as an error in every genre would violate
the product boundary.

textlint offers a pluggable natural-language linting model and machine-readable
formatters \citep{textlint2026docs}. It is useful for composing independent
checks and for feeding findings into an editor. Its Node ecosystem becomes a
maintenance consideration in a mostly Python and TeX repository; the adapter
should consume stable JSON rather than embed textlint's internal objects.

LanguageTool and LTeX+ provide grammar and language checks across ordinary
prose and, in the latter case, LaTeX-aware editing workflows
\citep{languagetool2026dev,ltexplus2026}. Their value is local correctness:
agreement, spelling, punctuation, and a limited class of usage problems. They
should be run after the argument and protected objects are stable. Otherwise
the author receives a large queue of low-consequence corrections before seeing
the missing mechanism.

proselint and Textidote occupy a similar but useful niche. They encode style
and usage advice, including vague language and passive constructions, and can
serve as additional rule sources. The system should retain the source rule and
version in each finding because a recommendation is not a fact about the
document. A former professor may deliberately choose a passive construction to
foreground a result, and a methods section may need a different register from
an executive opening.

When a sentence carries a citation, two different checks follow. First, a
resolver checks whether the key, identifier, title, and author record agree.
Second, a human or a checked source-support process asks whether the cited work
supports the sentence.
GROBID can help recover bibliographic fields from PDFs, but the current
inventory treats it as a candidate pending local deployment and fixture
results \citep{grobid2026}. No resolver should be named a ``citation validator''
if it checks identity only.

\section{AI-tell tools and the temptation to optimize the surface}
\label{sec:software-tells}

sloptrim and the related slop-forensics project are useful because they turn a
maintained pattern list into inspectable counts and spans. The repository audit
found a calibration in which a mechanically clean rejected draft scored better
than a document a human reader accepted. That result is not a failure of the
tools; it is a boundary observation. They can direct an author's attention to
repeated vocabulary, formulaic transitions, and sentence shapes. They cannot
decide whether the document has a mechanism or a persuasive reason to fund it.

fast-detect-gpt implements a probability-curvature family of AI detectors. Its
official repository and benchmark literature make it a useful reference for
aggregate drift experiments, but the false-positive and evasion results in
Chapter~\ref{sec:literature} rule out per-document use. Any adapter must
discard individual labels and retain only corpus-level distributions, with a
minimum cell size and a human interpretation of a change.

The same rule applies to the installed humanizer guidance. It supplies useful
patterns, such as inflated significance, vague sources, repeated headings, and
formulaic outlook sections. It is not a substitute for a writer's ear or for a
technical explanation. The main route therefore presents these projects as
diagnostic components; this appendix retains their operational details.

\section{Adjacent products an executive will compare}
\label{sec:software-adjacent}

An implementation proposal should not define its market only by the packages
its authors have already installed. A sponsor will reasonably ask why the team
cannot buy an existing writing assistant, use the review functions in its
LaTeX editor, or add a specialist academic service to the current workflow.
Use a deliberately modest standard to answer that question. A vendor's
documentation establishes a documented capability and a stated trust boundary.
It does not establish effectiveness on the Humanvoice population, and a
product page is not a substitute for a held-out comparison.

\subsection{General communication assistants}

Grammarly's current enterprise offering is broad: it works across the tools in
which teams write and combines writing suggestions with a more general agent
layer \citep{grammarly2026enterprise}. Its privacy documentation describes
how generative features may pass prompt text and context to service providers,
even while enterprise controls and consent are offered
\citep{grammarly2026privacy}. This is a credible alternative for ordinary
business communication. It can reduce the cost of local wording edits and may
be the correct incumbent for a short memorandum.

The documented strength is also the boundary. The assistant's unit is a
selected passage or a communication context. Humanvoice's unit is a versioned
technical document with equations, data, citations, claims, and a named
decision. A rewrite that is excellent at tone can still change a denominator,
drop a qualification, or make a claim sound more certain. The product wedge
therefore does not compete on generic fluency. It asks whether an assistant
can be placed behind a source map and a protected comparison, with the reader
outcome measured rather than inferred from acceptance of a suggestion.

\subsection{Academic writing assistants}

Writefull is closer to the target genre. Its Overleaf integration advertises
language support trained on published research and includes revision, style,
paraphrase, and other writing functions inside the LaTeX editor
\citep{writefull2026overleaf}. That makes it a serious candidate for a
language-diagnostic adapter or an incumbent arm. It does not, on the
documentation inspected, provide the complete Humanvoice record: a declared
reading brief, object-level correspondence for equations and claims, and a
reader decision that can be compared across workflows. Those absences are
questions for a fixture run, not accusations about the product.

Paperpal presents a broader academic suite, including contextual rewriting,
reference discovery, citation styles, submission checks, and integrations for
researchers \citep{paperpal2026assistant}. Its breadth is useful evidence that
the market already pays for language and submission support. It also creates a
measurement problem. Several checks and generative functions arrive in one
interface, so a positive user reaction cannot identify which component changed
the document or whether a protected object moved. Paperpal's stated privacy
position should be recorded in procurement, but the product still needs a
local or redacted path for restricted manuscripts.

Trinka is another relevant comparison because it positions itself for academic
and technical writing and documents privacy, institutional, integration, and
API options \citep{trinka2026privacy}. This may make it a plausible private
deployment or grammar adapter. The same discipline applies: a privacy feature
is not a proof of source correspondence, and an API is not a reader study.
The benchmark should test whether a Trinka finding can be tied to a source span,
whether an author can reject it without losing provenance, and whether its
processing mode meets the document owner's declared boundary.

\subsection{Authoring and review environments}

Overleaf already supplies several functions a new product should not
reimplement. Its documentation describes collaborative editing, comments,
version history, and Track Changes for accepting or rejecting collaborators'
edits \citep{overleaf2026trackchanges}. These features solve a real social
problem: multiple authors can inspect a manuscript and discuss a proposed
change in context. Humanvoice should integrate with that workflow where it is
the chosen incumbent, export a stable source snapshot, and return a reader
packet that does not expose private review history.

Track Changes remains a change display, not a semantic invariant. It tells a
reviewer which source text was edited and who made the edit; it does not by
itself establish that an equation is equivalent, a citation still supports a
claim, or a reader can make the intended decision. A product that treated
Overleaf's accept/reject control as proof of meaning preservation would repeat
the same category error as treating a grammar score as proof of comprehension.

\subsection{The comparison in one view}

The following table records the buy, adapt, or build question at the level of
the job rather than the brand. ``Documented'' means that the vendor describes
the capability. ``Unestablished'' means that the Humanvoice benchmark still
needs to observe it on a locked fixture and a named reader.

\begin{table}[H]
\centering
\small
\begin{tabularx}{\textwidth}{@{}>{\raggedright\arraybackslash}p{0.20\textwidth}
                              >{\raggedright\arraybackslash}p{0.24\textwidth}L
                              >{\raggedright\arraybackslash}p{0.23\textwidth}@{}}
\toprule
Alternative & Documented strength & Missing link for this proposal & Initial
disposition \\
\midrule
Grammarly / enterprise assistant & Broad inline suggestions, rewriting, and
organizational controls across common applications & No demonstrated
equation--claim lineage or named-reader endpoint in the inspected material &
Retain as an incumbent comparator; do not rebuild its general grammar layer \\
Writefull for Overleaf & Research-trained language support integrated with a
LaTeX editor & Source-level protected-object record and reader-outcome
measurement remain unestablished & Benchmark as a language adapter and
possible incumbent \\
Paperpal & Academic editing, references, submission checks, and multiple
integrations & One interface combines many interventions; causal contribution
and protected correspondence are not identified & Include in market scan;
test only declared functions \\
Trinka & Academic/technical language assistance with privacy and integration
options & Efficacy, source maps, and reader acceptance for the target genres
remain unestablished & Evaluate as a private-service or API candidate \\
Overleaf review layer & Comments, collaboration, version history, and
accept/reject Track Changes & Character/source changes are not an invariant
for equations, claims, or citations & Integrate as the authoring incumbent \\
Humanvoice core & Located findings, protected comparison, and a reader decision
by someone unfamiliar with the project in one record & No efficacy result and no broad market proof
yet & Build only the connective layer, then test its value \\
\bottomrule
\end{tabularx}
\caption{Adjacent products define the incumbent bundle and the falsifiable
Humanvoice wedge.}
\label{tab:software-adjacent}
\end{table}

\subsection{Buy, adapt, or build}

This comparison leads to a restrained implementation choice. Buy or reuse
general grammar, collaborative editing, and document-building capabilities
when their licensing and privacy terms fit the organization. Adapt a package
when it can emit a stable location, version, and explanation without taking
ownership of the document. Build the narrow connective layer that the current
alternatives leave implicit: a source snapshot, a protected-object schema, a
finding record, an author decision, and a reader outcome joined by immutable
identifiers.

That distinction affects the budget. Reimplementing a mature grammar engine
would spend the MVP budget on a capability that is already available. Reusing a
cloud editor without an object-level boundary would save engineering time while
moving the largest risk into an opaque service. The proposed pilot should price three configurations: a fully local stack; an
incumbent-editor stack with Humanvoice adapters; and a hosted academic
assistant used only on explicitly authorized excerpts. The measured reader and
operator burden, not a feature checklist, chooses among them.

\subsection{Tests that could overturn the position}

The product should be abandoned or repositioned if an adjacent product can
already produce the same source-to-reader record at lower total burden. The
test is concrete. Give each candidate the same source snapshot, protected
object manifest, and reading brief. Ask it to locate a deliberate
qualification loss, compare a changed number and citation, and produce a
reader packet that hides private workflow labels. Record abstentions, false
passes, operator minutes, network calls, and reader completion. A candidate
that wins this comparison should become the incumbent, even if it is not
open-source and even if it makes the Humanvoice implementation unnecessary.

Conversely, a candidate that wins a fluency preference but cannot explain a
protected change has not falsified the wedge; it has supplied a useful
component. This is why the software survey reports capability and evidence in
separate columns. The investment buys an option to learn whether the
connection among those components earns its cost, not a license to describe a
familiar product as a competitor it has not actually beaten.

\section{Local inference and reader evaluation}
\label{sec:software-inference}

llama.cpp is the leading local execution candidate in the inventory
\citep{llamacpp2026}. Its value is a deployment option when a document cannot
leave the organization's boundary. The product still needs a model card, a
prompt and sampling record, a resource profile, and a comparison against a
remote baseline. Local execution reduces one class of privacy exposure; it
does not make a generated suggestion correct or remove the need for disclosure.

Inspect AI supplies an evaluation harness for tasks, solvers, and graders
\citep{inspectai2026}. It can replay a fixed reader question, randomize version
order, and store machine-judge outputs beside human labels. It cannot recruit
the former professor whose decision defines the product. The first release
should keep it optional for model experiments.

The timed human task needs a simpler instrument. jsPsych can present local web
trials, randomize order, record responses and elapsed time, and export a
structured trial record \citep{jspsych2026}. A self-hosted jsPsych page is the
leading candidate for the feasibility reader packet because it does not make
model evaluation infrastructure a dependency of the human outcome. The
benchmark must still test keyboard accessibility, resume behavior, browser
logging, and export correctness; the library supplies instrumentation, not an
expert-reader sample.

\begin{table}[H]
\centering
\small
\begin{tabularx}{\textwidth}{@{}p{0.21\textwidth}p{0.24\textwidth}L p{0.22\textwidth}@{}}
\toprule
Candidate & First fixture & Disposition & Reason for
the boundary \\
\midrule
pylatexenc & Custom macro plus equation and table & adapt & Preserve raw spans;
abstain on unknown macros \\
unified-latex & Source positions, custom environments, and editor round trip &
pilot & Strong web-editor fit; unknown-construct behavior remains untested \\
TexSoup & Malformed input and compact Python traversal & pilot & Useful tolerant
fallback only if ambiguity remains explicit \\
plasTeX & Sections, labels, counters, and malformed environment & pilot &
Richer tree, but compatibility must match the canonical build \\
latexdiff & Number, citation, and equation moved across paragraphs & adapt &
Useful view, not semantic truth \\
Vale/textlint & Private label, vague referent, intentional technical term &
adopt as adapters & Findings need rule and context, not automatic rewrites \\
sloptrim & Accepted and rejected internal repair pair & pilot aggregate-only &
Calibrate against reader labels; no individual authorship score \\
llama.cpp & Same finding prompt under local and remote models & pilot & Measure
privacy, latency, and suggestion differences \\
Inspect AI & Order-swapped pairwise reader task & adapt & Organize replay; human
decision remains primary \\
jsPsych & Timed, randomized reconstruction task in a local browser & pilot &
Instrument human sessions without coupling them to a model harness \\
\bottomrule
\end{tabularx}
\caption{Initial software dispositions are hypotheses to be tested on a locked
fixture set.}
\label{tab:software-dispositions}
\end{table}

\section{The held-out fixture set}
\label{sec:software-fixtures}

The fixture set is deliberately small but varied. It contains a clean passage,
a passage with a private identifier, an equation whose exponent changes, a
table with a denominator mismatch, a citation with a valid identity but weak
support, a custom macro, a malformed environment, and a paragraph whose
repeated structure is intentional. Each fixture has a source file, a rendered
reference, a protected-object manifest, and an expected failure behavior.

For candidate $k$ and fixture $i$, record a capability result $C_{ki}$, a
false-pass indicator $S_{ki}$, runtime $t_{ki}$, and operator burden $b_{ki}$.
The benchmark summary is a matrix rather than a single package score:
\begin{equation}
  B_k = \{(C_{ki},S_{ki},t_{ki},b_{ki}): i=1,\ldots,I\}.
  \label{eq:software-benchmark-matrix}
\end{equation}
An adapter may be adopted only when it identifies the positive fixture, fails
or abstains on the negative fixture, and emits a location that a second
operator can replay. A fast false pass is worse than a slow abstention because
the former can change a scientific document without attracting attention.

The benchmark also records installation and maintenance facts. A package with
an incompatible license can remain a reference implementation. A package with
no recent release can be acceptable for a frozen migration and unsuitable for
a continuously operated service. A package that requires network access during
inference may be excluded from a private document even if its output is useful.
These are engineering facts, not aesthetic preferences, and they belong beside
the functional result.

\section{Composition and ownership}
\label{sec:software-composition}

The proposed stack has one owner for each boundary. The source adapter owns
locations and abstentions. The build wrapper owns the canonical PDF and
cross-reference diagnostics. The protected-object extractor owns normalized
equations, numbers, citations, and tables. The finding adapters own rule
identities and context. The local-model adapter owns prompts, sampling, and
network policy. The reader protocol owns the final outcome.

This division prevents a package from silently acquiring authority that its
evidence does not support. It also makes replacement possible. If a better
LaTeX parser appears, the source adapter can be replaced while the protected
object schema and reader protocol remain stable. If a model judge is retired,
the human labels and pairwise task remain. The product is the composition and
the evidence around it, not the brand of any one dependency.

The current inventory is therefore enough to begin a feasibility build, but not
enough to claim a validated software stack. The held-out benchmark, license
review, and independent replay remain release conditions. Part~III now turns
these dispositions into an architecture whose records can be implemented and
tested.
